% GENERATED FILE. Edit canonical lessons, then run npm run generate:books.
% canonical-source-hash: fnv1a64:dd287bfda422fbb7
% canonical-lessons: LA-C495-admiratio, LA-C495-miror, LA-C495-stupeo, LA-C495-attonitus, LA-C495-obstupesco

\chapter{Wonder and amazement}
\label{ch:la-wonder-and-amazement}
\hlchaptermodality{495}

\begin{chapteropening}
\textbf{By the end of this chapter:} \emph{I can describe wonder and amazement.}

Complete the last lesson of chapter 495: I can describe wonder and amazement.
\end{chapteropening}

\section[admīrātiō, admīrātiōnis]{admīrātiō, admīrātiōnis — admiration, wonder}

\label{lesson:LA-C495-admiratio}



\begin{quote}
Before the new one: say the Latin for disgust, then the Latin for it wearies.
\end{quote}

\subsection*{You'll want to know: admīrātiō, admīrātiōnis}
\textbf{admīrātiō, admīrātiōnis} — \textquotedblleft{}admiration, wonder\textquotedblright{}.

\begin{quote}
Hospitēs templum novum cum admīrātiōne spectant.
\end{quote}

\begin{quote}
\emph{The guests gaze at the new temple in wonder.}
\end{quote}

Feminine, third declension, from \emph{admīror}, already known. English \emph{admiration}.

\begin{cousinweb}
\textbf{admīrātiō} is the action noun of \emph{admīror} \textquotedblleft{}to wonder at\textquotedblright{}, built on \emph{ad-} and \emph{mīror} \textquotedblleft{}to wonder\textquotedblright{} (Lewis \& Short; EtymDB). English \emph{admiration} comes from it.
\end{cousinweb}

\begin{grammarlens}[title={What you've built}]
One more word for telling, explaining and describing at B1.
\end{grammarlens}

\subsection*{Your turn}
\begin{itemize}
\raggedright
  \item \emph{Say it:} \emph{admīrātiō, admīrātiōnis}
  \item \emph{Say it:} \emph{admīrātiō, admīrātiōnis}, once more
  \item \emph{Say it:} a sentence of your own with \emph{admīrātiō, admīrātiōnis}
  \item \emph{Recall:} say the Latin for disgust, then the Latin for it wearies, then say \emph{admīrātiō, admīrātiōnis} again
\end{itemize}

\begin{tcolorbox}[breakable,colback=teal!4,colframe=teal!35!black,title={Before you move on}]
What does admīrātiō, admīrātiōnis mean? (\textquotedblleft{}Admiration, wonder\textquotedblright{}.) Say it once more.
\end{tcolorbox}
\section[mīror, mīrārī]{mīror, mīrārī — to wonder, marvel at}

\label{lesson:LA-C495-miror}



\begin{quote}
Before the new one: say the Latin for it wearies, then the Latin for admiration.
\end{quote}

\subsection*{You'll want to know: mīror, mīrārī}
\textbf{mīror, mīrārī} — \textquotedblleft{}to wonder, marvel at\textquotedblright{}.

\begin{quote}
Puer elephantum mīrātur.
\end{quote}

\begin{quote}
\emph{The boy marvels at the elephant.}
\end{quote}

Deponent: \emph{mīror, mīrārī, mīrātus sum}. With the accusative, or with \emph{quod} — I am surprised that. English \emph{miracle} belongs to the same family.

\begin{cousinweb}
\textbf{mīror} is built on \emph{mīrus} \textquotedblleft{}wonderful\textquotedblright{}, which EtymDB traces through Proto-Italic to a Proto-Indo-European root meaning \textquotedblleft{}to smile\textquotedblright{}. English \emph{admire} and \emph{miracle} belong to the same family.
\end{cousinweb}

\begin{grammarlens}[title={What you've built}]
One more word for telling, explaining and describing at B1.
\end{grammarlens}

\subsection*{Your turn}
\begin{itemize}
\raggedright
  \item \emph{Say it:} \emph{mīror, mīrārī}
  \item \emph{Say it:} \emph{mīror, mīrārī}, once more
  \item \emph{Say it:} a sentence of your own with \emph{mīror, mīrārī}
  \item \emph{Recall:} say the Latin for it wearies, then the Latin for admiration, then say \emph{mīror, mīrārī} again
\end{itemize}

\begin{tcolorbox}[breakable,colback=teal!4,colframe=teal!35!black,title={Before you move on}]
What does mīror, mīrārī mean? (\textquotedblleft{}To wonder, marvel at\textquotedblright{}.) Say it once more.
\end{tcolorbox}
\section[stupeō, stupēre]{stupeō, stupēre — to be stunned, be amazed}

\label{lesson:LA-C495-stupeo}



\begin{quote}
Before the new one: say the Latin for admiration, then the Latin for to wonder.
\end{quote}

\subsection*{You'll want to know: stupeō, stupēre}
\textbf{stupeō, stupēre} — \textquotedblleft{}to be stunned, be amazed\textquotedblright{}.

\begin{quote}
Agricola in urbe magnā stupet.
\end{quote}

\begin{quote}
\emph{The farmer is stunned in the great city.}
\end{quote}

\emph{stupeō, stupēre, stupuī}. Struck dumb with wonder or fear. English \emph{stupor} comes from it, and \emph{stupid} from its adjective \emph{stupidus}.

\begin{cousinweb}
\textbf{stupeō} is inherited and usually traced to a Proto-Indo-European root meaning \textquotedblleft{}to strike, to stun\textquotedblright{} (EtymDB). English \emph{stupid} comes from its adjective \emph{stupidus}.
\end{cousinweb}

\begin{grammarlens}[title={What you've built}]
One more word for telling, explaining and describing at B1.
\end{grammarlens}

\subsection*{Your turn}
\begin{itemize}
\raggedright
  \item \emph{Say it:} \emph{stupeō, stupēre}
  \item \emph{Say it:} \emph{stupeō, stupēre}, once more
  \item \emph{Say it:} a sentence of your own with \emph{stupeō, stupēre}
  \item \emph{Recall:} say the Latin for admiration, then the Latin for to wonder, then say \emph{stupeō, stupēre} again
\end{itemize}

\begin{tcolorbox}[breakable,colback=teal!4,colframe=teal!35!black,title={Before you move on}]
What does stupeō, stupēre mean? (\textquotedblleft{}To be stunned, be amazed\textquotedblright{}.) Say it once more.
\end{tcolorbox}
\section[attonitus, attonita, attonitum]{attonitus, attonita, attonitum — thunderstruck, astonished}

\label{lesson:LA-C495-attonitus}



\begin{quote}
Before the new one: say the Latin for to wonder, then the Latin for to be stunned.
\end{quote}

\subsection*{You'll want to know: attonitus, attonita, attonitum}
\textbf{attonitus, attonita, attonitum} — \textquotedblleft{}thunderstruck, astonished\textquotedblright{}.

\begin{quote}
Attonitī nūntium audīmus.
\end{quote}

\begin{quote}
\emph{We hear the news, thunderstruck.}
\end{quote}

The past participle of \emph{attonō} — to strike with thunder.

\begin{cousinweb}
\textbf{attonitus} is the past participle of \emph{attonō} \textquotedblleft{}to strike with thunder\textquotedblright{}, built on \emph{ad-} and \emph{tonō} \textquotedblleft{}to thunder\textquotedblright{} (Lewis \& Short). English \emph{astonish} comes from a related compound of \emph{tonō}.
\end{cousinweb}

\begin{grammarlens}[title={What you've built}]
One more word for telling, explaining and describing at B1.
\end{grammarlens}

\subsection*{Your turn}
\begin{itemize}
\raggedright
  \item \emph{Say it:} \emph{attonitus, attonita, attonitum}
  \item \emph{Say it:} \emph{attonitus, attonita, attonitum}, once more
  \item \emph{Say it:} a sentence of your own with \emph{attonitus, attonita, attonitum}
  \item \emph{Recall:} say the Latin for to wonder, then the Latin for to be stunned, then say \emph{attonitus, attonita, attonitum} again
\end{itemize}

\begin{tcolorbox}[breakable,colback=teal!4,colframe=teal!35!black,title={Before you move on}]
What does attonitus, attonita, attonitum mean? (\textquotedblleft{}Thunderstruck, astonished\textquotedblright{}.) Say it once more.
\end{tcolorbox}
\section[obstupēscō, obstupēscere]{obstupēscō, obstupēscere — to be struck dumb, be astounded}

\label{lesson:LA-C495-obstupesco}



\begin{quote}
Before the new one: say the Latin for to be stunned, then the Latin for thunderstruck.
\end{quote}

\subsection*{You'll want to know: obstupēscō, obstupēscere}
\textbf{obstupēscō, obstupēscere} — \textquotedblleft{}to be struck dumb, be astounded\textquotedblright{}.

\begin{quote}
Hospitēs ad hanc fābulam obstupēscunt.
\end{quote}

\begin{quote}
\emph{The guests are astounded at this story.}
\end{quote}

\emph{obstupēscō, obstupēscere, obstupuī}. The inceptive of \emph{stupeō}: the moment the wonder strikes.

\begin{cousinweb}
\textbf{obstupēscō} is built on \emph{ob-} and \emph{stupēscō}, the inceptive form of \emph{stupeō} \textquotedblleft{}to be stunned\textquotedblright{} (Lewis \& Short; EtymDB).
\end{cousinweb}

\begin{grammarlens}[title={What you've built}]
One more word for telling, explaining and describing at B1.
\end{grammarlens}

\subsection*{Your turn}
\begin{itemize}
\raggedright
  \item \emph{Say it:} \emph{obstupēscō, obstupēscere}
  \item \emph{Say it:} \emph{obstupēscō, obstupēscere}, once more
  \item \emph{Say it:} a sentence of your own with \emph{obstupēscō, obstupēscere}
  \item \emph{Recall:} say the Latin for to be stunned, then the Latin for thunderstruck, then say \emph{obstupēscō, obstupēscere} again
\end{itemize}

\begin{tcolorbox}[breakable,colback=teal!4,colframe=teal!35!black,title={Before you move on}]
What does obstupēscō, obstupēscere mean? (\textquotedblleft{}To be struck dumb, be astounded\textquotedblright{}.) Say it once more.
\end{tcolorbox}
